% IW49091C
% Nadelstichverletzung
% 2.0.1
% 2015-11-15
\label{m:nadelstichverletzung}
\Z{Basiernd auf \cite{OeAG2007}.}
\begin{itemize}
    \item Erstmaßnahmen durchführen: \nocite{OeAG2007}
    \begin{itemize}
        \item bei Stich- bzw.Schnittverletzungen:
        Blutung durch Kompression des umgebenden
        Gewebes induzieren. Danach desinfizieren und einen
        in Desinfektionsmittel getränkten Tupfer 10\Min*
        auf der Wunde belassen.
        \item bei Spritzern mit Blut oder Sekreten in Auge oder Mund:
        Schleimhaut mit Wasser u./o. Schleimhautdesinfektionsmittel spülen
    \end{itemize}
    \item Leitstelle informieren
    \item ${\rightarrow}$ \EE{Umgehend in ein
    geeignetes Krankenhaus} (Von der Leitstelle erfragen).
    \item Patient: Zielspital informieren:
    \begin{itemize}
        \item Blut sichern für \Abk{HIV}- und Hepatitis-Serologie
    \end{itemize}
    \item  \Ca{Unfallmeldung}  (wichtig!), Blutabnahme
    \item ggf. \Abk{HIV}-Postexpositionsprophylaxe (\Abk{PEP})
    \begin{itemize}
        \item kurzfristige Therapie (4 Wochen)
        \item Einzelfall muss immer individuell beurteilt werden
        (Nutzen/Risiko), starke Nebenwirkungen.
        \item Muss innerhalb von 2\Hour* begonnen werden ${\rightarrow}$ daher
        \underline{\EE{sofort}} in ein geeignetes Krankenhaus!
    \end{itemize}
    \item Meldung an Betriebsleitung
    \item \EE{Unfallmeldung (wichtig!) an die \Abk{AUVA}}
\end{itemize}

\begin{Danger}
  \item Unfallmeldung nicht vergessen!
  
  Sofort in ein geeignetes Krankenhaus!
\end{Danger}
